\documentclass[10pt]{article}
\usepackage[margin=1in]{geometry}
\usepackage{amsmath,amssymb,amsthm,mathtools}
\usepackage[hidelinks]{hyperref}
\newtheorem{theorem}{Theorem}
\newtheorem{lemma}[theorem]{Lemma}
\newcommand{\E}{\mathbb E}
\newcommand{\Be}{\operatorname{Beta}}
\title{Real-shift coefficientwise log-concavity\\for rising-factorial ratios\\
\large A proof of MathDB problem 324472}
\author{Principia Math In Progress}
\date{September 12, 2026}
\begin{document}
\maketitle

\noindent\textbf{Task 1: solve the original problem.}
We prove Conjecture 1 of Kalmykov--Karp \cite[p.~408]{KK} under the
hypotheses of their Theorem 1, with an additional relaxation when the
parameter gap is at most one.

\begin{theorem}\label{main}
Let \(c\ge a>0\), and let \(f_j\) be a nonnegative log-concave sequence
without internal zeros. Define
\[
 F(\mu;x)=\sum_{j\ge0}f_j\frac{(a+\mu)_j}{(c+\mu)_j}\frac{x^j}{j!},
 \qquad (b)_j=\prod_{i=0}^{j-1}(b+i).
\]
For every \(\mu,\nu\ge0\), every coefficient of
\[
 F(\mu;x)F(\nu;x)-F(0;x)F(\mu+\nu;x)                 \tag{1}
\]
is nonnegative. If \(0\le c-a\le1\), the same conclusion holds for every
nonnegative sequence \(f_j\), without the log-concavity restriction.
These are formal-series assertions; they imply the corresponding
Wright log-concavity for nonnegative \(x\) in the common convergence domain.
\end{theorem}

The proof uses elementary comparisons of densities. We recall the
one-crossing argument: if two probability densities on an interval have
an increasing ratio, then the second law is stochastically larger,
provided the ratio is second divided by first. Indeed the difference of
the normalized densities crosses zero at most once, from negative to
positive, so their distribution functions have the required order.
The same argument applies if the second support also extends farther
to the right and has no extra portion to the left. Decreasing ratios
reverse the conclusion.

\begin{lemma}\label{poly}
For symmetric nonnegative weights satisfying
\[
 w_k=w_{n-k},\qquad
 w_0\le w_1\le\cdots\le w_{\lfloor n/2\rfloor},      \tag{2}
\]
the Bernstein polynomial
\[
 R_w(r)=\sum_{k=0}^n\binom nk w_k r^k(1-r)^{n-k}
 \tag{3}
\]
is nondecreasing on \([0,1/2]\) and satisfies \(R_w(r)=R_w(1-r)\).
\end{lemma}
\begin{proof}
Symmetry is immediate. For \(n\ge1\), differentiate (3):
\[
 R_w'(r)=n\sum_{k=0}^{n-1}(w_{k+1}-w_k)
                  \binom{n-1}k r^k(1-r)^{n-1-k}.
 \tag{4}
\]
The coefficient differences at reflected indices are opposite, with the
lower-index one nonnegative. On \([0,1/2]\), the lower-index Bernstein
term is at least its reflection. Pairing terms proves (4) nonnegative.
The cases \(n=0,1\) are constant.
\end{proof}

\begin{lemma}\label{beta}
Let \(X\sim\Be(u,d)\) and \(Y\sim\Be(v,d)\) be independent, with
\(u,v,d>0\), and put
\[
 P_w(x,y)=\sum_{k=0}^n\binom nk w_kx^ky^{n-k}.
 \tag{5}
\]
At fixed \(S=u+v\), the expectation \(\E P_w(X,Y)\) is nondecreasing
as \(u,v\) are balanced, in either of the following cases:
\begin{enumerate}
\item \(0<d\le1\) and \(w\) is any symmetric nonnegative vector;
\item \(d>1\) and \(w\) satisfies (2).
\end{enumerate}
\end{lemma}
\begin{proof}
Write
\[
 q=\frac{\min(X,Y)}{\max(X,Y)},\qquad z=\max(X,Y),\qquad
 \delta=\frac{|u-v|}{2}.
\]
Both \(q,z\) lie in \((0,1)\) with probability one. Adding the densities
from the two possible orders of \(X,Y\), with Jacobian \(z\), gives
the joint density
\[
 \frac{2}{B(u,d)B(v,d)}
 z^{S-1}q^{S/2-1}(1-z)^{d-1}(1-qz)^{d-1}
 \cosh(\delta\log q).                              \tag{6}
\]
In particular the conditional law of \(z\) given \(q\) does not depend
on \(\delta\). The marginal density of \(q\), apart from its normalizing
constant, has the form
\[
 A(q)q^{S/2-1}\cosh(\delta\log q),\qquad
 A(q)=\int_0^1 z^{S-1}(1-z)^{d-1}(1-qz)^{d-1}\,dz.  \tag{7}
\]
The integral is positive and finite for \(0<q<1\).
If \(\delta_2\ge\delta_1\ge0\), the ratio of the marginal densities
is decreasing in \(q\). To check this, put \(t=-\log q>0\) and note
\[
 \frac{d}{dt}\log\frac{\cosh(\delta_2t)}{\cosh(\delta_1t)}
 =\delta_2\tanh(\delta_2t)-\delta_1\tanh(\delta_1t)\ge0.
 \tag{8}
\]
Thus \(q\) is stochastically larger as \(\delta\) decreases.
Since \(P_w\) is symmetric, it remains to show that the
\(\delta\)-independent conditional mean
\[
 m(q)=\E[P_w(X,Y)\mid q]
 \tag{9}
\]
is nondecreasing in \(q\).

Suppose first that \(0<d\le1\). For \(0<q_1<q_2<1\), the ratio
of the conditional \(z\)-densities, apart from a positive constant, is
\[
 \left(\frac{1-q_2z}{1-q_1z}\right)^{d-1}.           \tag{10}
\]
The inner ratio decreases in \(z\), so (10) increases in \(z\).
Consequently the conditional moment \(\E[z^n\mid q]\) is nondecreasing
in \(q\). Also,
\[
 P_w(z,qz)=z^nQ_w(q),\qquad
 Q_w(q)=\sum_{k=0}^n\binom nk w_kq^{n-k}.
 \tag{11}
\]
The polynomial \(Q_w\) has nonnegative coefficients, hence is nonnegative
and nondecreasing on \((0,1)\). Formula (11) proves the required
monotonicity of \(m(q)\).

Now suppose \(d>1\), and set \(T=X+Y=(1+q)z\). By (6), its conditional
density given \(q\), apart from a positive normalizing constant, is
\[
 t^{S-1}\bigl(1-t+h(q)t^2\bigr)^{d-1},
 \quad 0<t<1+q,\qquad h(q)=\frac{q}{(1+q)^2}.        \tag{12}
\]
The quadratic is the product
\((1-t/(1+q))(1-qt/(1+q))\), so it is positive on that support.
Moreover \(h\) is increasing for \(0<q<1\). If \(q_2>q_1\), write
\(h_i=h(q_i)\). On the common support the ratio of the two quadratics is
\[
 \frac{1-t+h_2t^2}{1-t+h_1t^2}
 =1+(h_2-h_1)\frac{t^2}{1-t+h_1t^2}.                \tag{13}
\]
It is increasing with \(t\), since
\[
 \frac{d}{dt}\frac{t^2}{1-t+h_1t^2}
 =\frac{t(2-t)}{(1-t+h_1t^2)^2}>0
 \quad(0<t<1+q_1<2).                              \tag{14}
\]
Raising (13) to \(d-1>0\) preserves this monotonicity, and the upper
support endpoint \(1+q\) increases. The one-crossing argument therefore
shows that \(T\), conditionally on \(q\), is stochastically increasing
with \(q\). In particular \(\E[T^n\mid q]\) increases.

Homogeneity, symmetry, and Lemma~\ref{poly} give
\[
 P_w(X,Y)=T^nR_w\!\left(\frac{q}{1+q}\right).
 \tag{15}
\]
The second factor is nonnegative and nondecreasing in \(q\), since
\(q/(1+q)\in(0,1/2)\). Multiplying it by
\(\E[T^n\mid q]\) proves that (9) is nondecreasing. In both cases,
the stochastic comparison following (8) now proves the lemma.

All conditional-density arguments take place on \(0<q<1\). Possible
limiting singularities at \(q=1\) cause no exception: that endpoint has
probability zero, and the polynomial \(P_w\) is bounded on \([0,1]^2\).
\end{proof}

\begin{proof}[Proof of Theorem~\ref{main}]
If \(d=c-a=0\), every rising-factorial ratio is one and (1) vanishes.
Assume \(d>0\). For \(B_{u,d}\sim\Be(u,d)\), the beta integral gives
\[
 \E B_{u,d}^{\,k}=\frac{B(u+k,d)}{B(u,d)}
                =\frac{(u)_k}{(u+d)_k}.             \tag{16}
\]
For each degree \(n\), take \(w_k=f_kf_{n-k}\) in (5). The Cauchy
product and (16) identify \(n!\) times the coefficient of \(x^n\) in
(1) with
\[
 \E P_w(B_{a+\mu,d},B_{a+\nu,d})
 -\E P_w(B_{a,d},B_{a+\mu+\nu,d}),                  \tag{17}
\]
where the two beta variables in each expectation are independent.
The first parameter pair is at least as balanced as the second,
because their sums agree and \(|\mu-\nu|\le\mu+\nu\).

If \(f\) is log-concave without internal zeros, its product weights
satisfy (2). On the positive support this follows from
\[
 \frac{w_{k+1}}{w_k}
 =\frac{f_{k+1}/f_k}{f_{n-k}/f_{n-k-1}}\ge1
 \quad(k<n/2).                                    \tag{18}
\]
Its positive support is a centered interval, so initial or terminal
zeros cause no exception. If all weights vanish, the coefficient is zero.
Lemma~\ref{beta} proves (17) nonnegative. Its first case also proves
the stated extension for arbitrary nonnegative \(f\) when \(d\le1\).
This covers degree zero and vanishing shifts as well. Replacing
\((a,c)\) by \((a+t,c+t)\), \(t\ge0\), supplies every common base
shift required in the Wright definition.
\end{proof}

\noindent\textbf{Verification.}
The companion program passes 441 exact rational coefficient cases,
3 exact density-derivative identities, 225 conditional-moment
comparisons, and 3 independent transformed beta integrals. The analytic
proof, not these finite checks, establishes all degrees and all parameters.
Source and current claim checks are recorded in the literature audit.

\paragraph{Task 2.}
Determine when arbitrary nonnegative input can produce a negative
coefficient for \(c-a>1\), and find the smallest failing degree
as a function of the parameter gap and the base parameter.

\paragraph{Task 3.}
Determine all equality cases coefficient by coefficient and derive
sharp quantitative lower bounds for (1) for arbitrary positive real shifts.

\paragraph{Task 4.}
Extend the conditional-density comparison to products of three or more
beta variables and identify multivariate coefficient conditions that
give monotonicity under majorization of the shape parameters.

\paragraph{Task 5.}
Prove or refute analogous real-shift coefficient inequalities for
series containing several rising-factorial ratios, using joint
ratio conditioning to handle the additional parameter differences.

\begin{thebibliography}{9}
\bibitem{KK}
S.~I. Kalmykov and D.~B. Karp,
Log-convexity and log-concavity for series in gamma ratios and applications,
\emph{J. Math. Anal. Appl.} \textbf{406} (2013), 400--418.
\href{https://doi.org/10.1016/j.jmaa.2013.04.061}{doi:10.1016/j.jmaa.2013.04.061};
\href{https://arxiv.org/abs/1211.2882}{arXiv:1211.2882}.
\end{thebibliography}
\end{document}
